\chapter{Introduction}
\label{chap:introduction}

\section{Context and motivation}

LLM-based agents now execute actions without a human reviewing the decision beforehand: they refund customers, commit code to production repositories, and triage medical queries autonomously. These systems reason over multiple steps, invoke external tools, and pursue goals across long interaction horizons. As agents take on greater responsibility, an undetected failure has consequences ranging from a degraded user experience to an incorrect action taken on a user's behalf.

Evaluating these systems is difficult because agents break the assumptions of traditional evaluation methods, which judge correctness at the level of a single response. Grading a single response is already hard, since outputs are sampled from a stochastic policy rather than read from a deterministic one \cite{songGoodBadGreedy2024}, and open-ended tasks have no fixed answer to compare an output against \cite{liuGEvalNLGEvaluation2023}. Agents add a further difficulty: they produce an execution \emph{trajectory}, the ordered sequence of reasoning steps, tool calls, and intermediate outputs that the agent goes through before it returns a final response. This trajectory interacts with live external systems, and it exposes failures that a single response cannot show, including tool selection errors, reasoning mistakes, and coordination breakdowns across agents \cite{maDemystifyingLifecycleFailures2026, mohammadiEvaluationBenchmarkingLLM2025}. A single root-cause error frequently cascades through subsequent decisions and amplifies its impact far beyond the step where it originated \cite{zhuWhereLLMAgents2025}.

Methods that inspect the trajectory exist, but they remain scattered across the categories of evaluation techniques \cite{mohammadiEvaluationBenchmarkingLLM2025}. A fixed benchmark that is run once before release does not capture how an agent's behavior evolves after deployment, because system-level interactions continue to shape it over time. Xia et al. review academic and industrial practice and propose to treat evaluation as a continuous activity spanning development and production time, rather than as a check performed once, before release. In their proposal, the evidence that evaluation produces drives both adaptation in production and later redevelopment \cite{xiaEvaluationDrivenDevelopmentOperations2025}.

The unit of analysis is therefore the trajectory, not the final response. A metric that scores only the output cannot distinguish a run that reaches the correct answer through a sound chain of tool calls from a run that reaches the same answer through an ungrounded guess. In a multi-agent system the unit of analysis extends to the trajectory of every participating agent, together with the messages exchanged between them: what an agent sends and the actions it takes become the input of the agents that run after it \cite{guoLargeLanguageModel2024}. Evaluation must therefore target how the agents communicate and delegate work to each other, not any agent in isolation \cite{zhuLLMBasedMultiAgentOrchestration2026a}.

Most multi-agent systems arrange their agents into a fixed structure, referred to as a \emph{topology}, that fixes a set of interaction paths between agents. In some systems the topology may also change at runtime: agents can be added, removed, or reconnected as the task unfolds \cite{zhuLLMBasedMultiAgentOrchestration2026a}. Among the paths a topology fixes, one is selected for each input: this choice determines which agent receives a given unit of work, and it is itself made by an agent. This decision is referred to here as \emph{routing}, and the component that takes it as the \emph{coordinator}. The coordinator decomposes the incoming task, assigns subtasks to worker agents, collects their outputs, and decides what happens next; typically, a language model reads the current state to select the next worker \cite{zhuLLMBasedMultiAgentOrchestration2026a}. In a hierarchical system, a coordinator takes the same decision at every level of the tree \cite{zhuLLMBasedMultiAgentOrchestration2026a}. Zhu et al. survey multi-agent orchestration from 2023 to early 2026 and list hierarchical coordination among the dominant patterns in six of the seven application domains they review, including software development, customer service, and scientific research \cite{zhuLLMBasedMultiAgentOrchestration2026a}. The same survey recommends the topology for large decomposable tasks with clear boundaries between domains, and for systems of more than 10 to 15 agents \cite{zhuLLMBasedMultiAgentOrchestration2026a}. Tooling supports the arrangement directly: provider agent SDKs expose handoffs between agents as first-class constructs \cite{OpenAIAgentsSDK}, and delegation to specialized sub-agents is catalogued as a general design pattern \cite[Chap.~7]{gulliAgenticDesignPatterns2025}.

The routing decision has the form of a classification over natural language, taken from a set of alternatives declared in the coordinator's prompt. Sub-agents may be distinguished by topic, or they may be equally capable over the same subject matter and differ only in what they optimize for. Two requests about the same system could therefore require different sub-agents, despite the vocabulary of the input they receive being largely overlapping.

An error in this decision can determine a failure in the outcome of the rest of the run. The selected sub-agent might apply its own criteria correctly to a request that should have been handled under different criteria. Its output can therefore be internally consistent and still address the wrong problem, and every step that follows the delegation builds on it \cite{zhuWhereLLMAgents2025}. Failures often originate in the coordination between agents rather than inside any one of them. In an annotated study of 1642 multi-agent traces, Cemri et al. attribute 32.3\% of failures to this category, which they label inter-agent misalignment \cite{cemriWhyMultiAgentLLM2025a}. Routing is the decision that establishes this interaction: it determines which agent receives the input, and therefore what the rest of the run builds on.

The outcome of a run shows only that the run failed, not that a routing decision caused the failure. Assigning an observed failure to the step that produced it is the problem of \emph{failure attribution}. On the Who\&When dataset of annotated failure logs, existing methods identify neither the responsible agent nor the decisive step reliably \cite{zhangWhichAgentCauses2025}. Establishing that a given run was misrouted therefore requires a human to inspect what the coordinator did, and this does not scale to the volume of runs that a deployed system produces. Checking the routing decision when the coordinator takes it, before the selected sub-agent runs, avoids this dependence on post-hoc attribution, and it also allows the run to be stopped before the selected sub-agent starts running. A signal available at that point can hold the request for review or select a different sub-agent, before the tokens and tool calls of a misrouted run are spent.

Before the model emits the sequence of tokens that expresses its routing choice, it computes a hidden state at every layer. This state exists before the decision becomes visible in the output, and it is not limited to what the model later states. \emph{Probing} is the technique of training a classifier on a model's hidden states to recover such a property, and prior results show that properties of interest can be recovered this way. Wu et al. probe a model's activation space for the tool that it is about to call, and find that the choice is encoded along a linear direction, readable before the call is generated \cite{wuToolCallingLinearly2026}. Azaria and Mitchell probe for a property that the model is never asked to report at all: they recover whether a statement is true from the activations of a model that is not asked to judge it \cite{azariaInternalStateLLM2023}. Both results follow the same pattern: a property probed from the hidden state can precede, and differ from, what the model eventually states.

Similar work has already applied probing to a routing decision. Ratnakar et al. train a linear probe on the hidden states of a frozen small language model to classify a query directly against a pool of 100 specialized agents, to reduce the latency and cost of a prompt-based LLM router \cite{ratnakarLatentGateLowLatencySemantic2026}. The frozen model is never instructed to route, so the probe is never compared against the model's own routing accuracy. Ratnakar et al. validate the probe on CLINC150, a separate 150-class intent-classification benchmark, where it outperforms a bag-of-words classifier, an SVM trained on TF-IDF features, by a small margin, 92.9\% against 91.2\% \cite{ratnakarLatentGateLowLatencySemantic2026}. This small margin suggests that most of the routing signal on CLINC150 is already recoverable from the query text alone. On their own 100-agent benchmark, they compare the probe only against a sentence-embedding classifier, which the probe outperforms by a wider margin, 80.0\% against 60.0\% \cite{ratnakarLatentGateLowLatencySemantic2026}. The evaluation is also limited to single-turn queries, and the pool of 100 candidate agents is far larger than the roughly three to seven a production coordinator typically selects among \cite{ratnakarLatentGateLowLatencySemantic2026}. It remains an open problem whether a hidden-state probe is more accurate than the coordinator itself when the candidate pool is small and sub-agents share overlapping terminology. In that setting, a lexical classifier cannot reliably resolve the correct sub-agent, so such a probe could serve as a check on the coordinator's own routing decision. It is also an open problem whether the same holds in a multi-turn setting.

\section{Research questions}

This work addresses whether the routing decision can be checked at the time it is taken, from the internal state of the model that takes it, on requests where a classifier trained on the query text alone fails to recover the correct sub-agent. Two interaction settings are used: in the single-turn setting, the coordinator routes one isolated request; in the multi-turn setting, it routes the last request of a conversation, so the state at the last prompt token also carries the preceding turns.
\phantomsection\label{rq:intro-rq1}
\begin{rqbox}[RQ1]
On requests where the correct sub-agent is difficult to determine from the wording alone, does a probe on the hidden states of a coordinator model recover the correct sub-agent more accurately than the coordinator's own selection?
\end{rqbox}


\phantomsection\label{rq:intro-rq2}
\begin{rqbox}[RQ2]
Does a probe on the coordinator's hidden states recover the correct sub-agent more accurately than classifiers trained on a fixed representation of the query text, such as raw word counts, TF-IDF weighting, or sentence embeddings?
\end{rqbox}


\phantomsection\label{rq:intro-rq3}
\begin{rqbox}[RQ3]
Does the optimal probe layer fall at a consistent relative depth across different LLMs?
\end{rqbox}


\phantomsection\label{rq:intro-rq4}
\begin{rqbox}[RQ4]
Can the disagreement between the probe and the coordinator be used as a runtime flag for likely misrouting?
\end{rqbox}


\section{Approach}

Two synthetic datasets of software developer requests are constructed to answer the research questions, one single-turn and one multi-turn. Both are organized around five sub-agents: time-to-market, scalable architecture, resource optimization, legacy compatibility, and strict regulatory compliance. Each request describes a technical scenario built around one or more \emph{constraints}: technical requirements the system must satisfy, such as a delivery deadline or a data-privacy rule. In each request, one constraint is \emph{binding}: it is imposed by someone with authority, such as a manager or a customer, and it determines the correct sub-agent. The other constraints are \emph{waived}: they are explicitly set aside, such as moved to a later milestone, even if they remain part of the request. For instance, a request may state a hard delivery deadline and also state that the system does not need to meet the compliance standards required for production; here, the deadline is binding, and compliance is waived. Requests are generated at three levels of difficulty:

\begin{itemize}
    \item \emph{Plain}: only one sub-agent can satisfy the request, and no other constraint is mentioned.
    \item \emph{Natural}: several constraints are mentioned, but the one that is binding is stated plainly.
    \item \emph{Adversarial}: two requests share almost the same wording and differ only in the detail that decides which sub-agent is correct, so one request acts as a decoy for the other.
\end{itemize}

Adversarial requests are constructed to be difficult to classify on wording alone, since the two requests in a pair use language that is close between sub-agents. The multi-turn dataset follows the same categories and the same structure, and distributes the elements of a request over a multi-turn conversation.

Eight open-weight models are run as coordinators on both datasets; for each request, the selected sub-agent and the hidden states at the last prompt token are recorded at every layer. A linear probe is then trained per layer to predict the correct sub-agent from those hidden states. Its accuracy is compared against the accuracy of the coordinator itself and against text-only baselines computed on the query. The disagreement between the probe and the coordinator is then evaluated as a runtime flag for likely misrouting.

\section{Contributions}

The work's contribution is two datasets for lexically hard routing decisions, one single-turn and one multi-turn, sharing the same sub-agents and scenario pool so the two interaction settings can be compared directly. To establish that the datasets are a hard test of routing, text-only classifiers measure how much of the label is recoverable from a fixed representation of the request text.

The two datasets and the source code are publicly available at \url{https://github.com/canta2899/misroute}. The release tagged \texttt{v1.0.0} corresponds to the version described in this work.

\section{Structure of the work}

This work is organized in nine chapters.

Chapter~\ref{chap:llm} covers how large language models work, from tokenization to the final output distribution, and introduces the concept of hidden states, the intermediate representations a model computes at each layer and token position.

Chapter~\ref{chap:agents} then describes how large language models can be contextualized as the decision-making component of agent systems. It introduces the mechanisms that let an LLM-based agent perceive, decide, and act, and the concept of a multi-agent system together with its common topologies. Among these topologies, the coordinator pattern is the one the rest of the work addresses.

Chapter~\ref{chap:evaluating-agents} reviews how large language models and LLM-based agents are evaluated, from reference-based metrics and benchmarks to holistic evaluation, LLM-as-a-judge, and trajectory-aware evaluation for multi-agent systems. This leads to failure attribution, the problem of tracing an observed failure back to the step that caused it. The chapter applies this problem to the coordinator's routing decision, the target that the rest of the work investigates.

Chapter~\ref{chap:motivation} compares candidate designs for evaluating a coordinator's routing decision, and reviews related literature on hidden-state probing. It then identifies the four research questions this work targets, and the requirements for the rest of the work.

Chapter~\ref{chap:dataset-construction} addresses the construction of two synthetic datasets (single-turn and multi-turn) according to the requirements established in Chapter~\ref{chap:motivation}.

Chapter~\ref{chap:probing-results} describes the experimental setup and reports the results obtained across eight open-weight models. Chapter~\ref{chap:results-analysis} analyzes those results and answers the research questions.

Finally, Chapter~\ref{chap:conclusions} identifies the limitations of the work and outlines directions for future work.
